\section{Bandit alerts and triage}
\label{sec:triage}

\Cref{sec:discovery} named the weaknesses that HW1's alerts point to and the attacker against whom
each one matters. This section describes the alerts themselves, which are what the agent you build in
HW1 receives and must label. We show how Bandit reads a program, what its rules and alert records
contain, and when a matched weakness is a vulnerability, using alerts from the development set with
and without a real flaw behind them. We then explain why every alert needs triage, what each kind of
triage mistake costs, how alerts are prioritized, and what a label must carry.

\subsection{How Bandit reads a program}
\label{sec:triage-tree}

Bandit finds alerts by matching shapes in the structure of a program, without running it. A
\textbf{syntax tree} is the tree that a parser builds from source code. Each node stands for one
construct of the language, such as an assignment, a function call or a name, and the children of a
node are the parts of that construct. Bandit parses each Python file into the syntax tree that
Python's own parser produces \citep{python-ast} and runs every rule against the nodes of the tree
\citep{bandit}.

A rule, as \cref{sec:discovery} defined it, is a pattern, a message and metadata. In Bandit, the
pattern is a shape in the syntax tree, for example a call to a function with a given name. When a
node has that shape, the rule reports an alert at the line where the node occurs, together with the
rule's message and its severity (\cref{sec:triage-severity}).

What a rule can see is limited to the part of the tree it reads. A rule reads the node it matched and
a few nodes near it, which record the call as it is written: which function is called and what is
passed to it. The rule does not follow a name back through the assignments that gave the name its
value. It therefore learns nothing about where an argument's value came from. When the argument is a
variable, as \texttt{fb} is in \cref{fig:triage-tree}, the rule cannot tell whether the variable holds
a value read from a request or a constant assigned a few lines earlier.

Every Bandit alert reports a matched shape and leaves open where the data at that shape comes from.
Whether data from an attacker can reach the matched call is the question that triage answers
(\cref{sec:triage-undecided}).

\begin{figure*}[t]
  \centering
  \begin{adjustbox}{width=\textwidth}% One Bandit match, in three panels: one line of Radicale, its syntax tree, and the name list of
% rule B301. The highlighted path is what the rule reads: the called name, pickle.load. The
% argument fb is a leaf the rule never visits. Redrawn for the notes from
% lecture-3--2026-09-25/slides/figures/hw1-ast-rule.tikz (lecture3-hw.pptx slide 21), without the
% file-path labels. Styles are local so the figure needs nothing beyond the notes' preamble.
\begin{tikzpicture}[
  tn/.style={gnode, font=\sffamily\small, minimum height=7mm, inner sep=4pt},
  card/.style={draw=ibmfg!60, rounded corners=3pt, fill=white, inner sep=6pt, align=left,
               font=\sffamily\small},
  hl/.style={draw=ibmblue, line width=1.3pt},
  tree/.style={draw=ibmfg!45, thick},
  path/.style={draw=ibmblue, line width=1.3pt},
  hit/.style={-{Stealth[length=2mm]}, draw=ibmblue, line width=1.2pt, rounded corners=3pt},
  arr/.style={-{Stealth[length=2mm]}, thick, draw=ibmfg!80, rounded corners=3pt},
  figcap/.style={font=\sffamily\small, inner sep=1pt, align=center, text=black!60},
  tag/.style={font=\sffamily\small, fill=white, inner sep=1pt, align=center},
]
  % left: the line of code
  \node[card, font=\ttfamily\small, minimum height=10mm] (code) at (0,0.6)
        {(version, row) = pickle.load(fb)};
  \node[figcap, below=1.5mm of code] {one line of Radicale};

  % middle: the tree
  \node[tn] (assign) at (7.0, 2.4) {Assign};
  \node[tn] (tuple) at (5.2, 1.1) {Tuple};
  \node[tn, hl] (call) at (8.8, 1.1) {Call};
  \node[tn] (v) at (4.0,-0.4) {Name\,version};
  \node[tn] (r) at (6.2,-0.4) {Name\,row};
  \node[tn, hl] (attr) at (8.2,-0.4) {Attribute};
  \node[tn] (fb) at (10.2,-0.4) {Name\,fb};
  \node[tn, hl] (pk) at (7.1,-1.9) {Name\,pickle};
  \node[tn, hl] (ld) at (9.2,-1.9) {attr\,load};
  \draw[tree] (assign) -- (tuple);
  \draw[path] (assign) -- (call);
  \draw[tree] (tuple) -- (v);
  \draw[tree] (tuple) -- (r);
  \draw[path] (call) -- (attr);
  \draw[tree] (call) -- (fb);
  \draw[path] (attr) -- (pk);
  \draw[path] (attr) -- (ld);
  \node[figcap] at (7.0,-2.9) {its syntax tree};

  % right: the rule
  \node[card, text width=33mm] (rule) at (13.9, 0.6)
        {\textbf{Rule B301}\\[2pt]
         names to match:\\
         {\ttfamily\small pickle.loads}\\
         {\ttfamily\small\color{ibmblue}\bfseries pickle.load}\\
         {\ttfamily\small pickle.Unpickler}\\
         {\ttfamily\small dill.loads}\ \ldots};
  \node[figcap, below=1.5mm of rule] {a list of names};
  \node[tag, below=6mm of rule, text=ibmblue, font=\sffamily\small\bfseries]
        {match: alert B301 on this line};

  % the arrows between panels
  \draw[arr] (code.east) -- ++(0.7,0) |- node[tag, pos=0.75, above] {parse} (assign.west);
  % past Name fb, never through it: along the leaf row, then up, then into the card
  \coordinate (entry) at ($(rule.west)+(0,-0.15)$);
  \coordinate (turn) at ($(fb.east)+(0.45,0)$);
  \draw[hit] (ld.east) -- (ld.east -| turn) -- (turn |- entry) -- (entry);
  \node[tag, anchor=south west, text=ibmblue] at ($(ld.east)+(0.12,0.1)$) {pickle.load};
\end{tikzpicture}
\end{adjustbox}
  \caption{One line of Radicale, its syntax tree and Bandit's rule B301. The highlighted path leads
  from the assignment to the called name, \texttt{pickle.load}, which the rule finds in its list of
  names and reports as an alert. The argument \texttt{fb} is a leaf that the rule never reads.}
  \label{fig:triage-tree}
\end{figure*}

\Cref{fig:triage-tree} follows one match from a line of code to an alert. The line comes from
Radicale, a calendar server and one of HW1's held-out targets. It loads a stored object from the
open file \texttt{fb} and unpacks it into two variables. Python's parser turns the line into an
assignment node. Its target is a tuple of two names, and its value is a call. The function called is
the attribute \texttt{load} of the name \texttt{pickle}, and the one argument is the name
\texttt{fb}.

Bandit reads the called name, \texttt{pickle.load}, along the highlighted path and looks it up in
the list of names held by rule B301. The name is on the list, so B301 reports an alert on this line.
The argument appears in the tree only as the leaf \texttt{Name fb}, and the rule never visits it.
Even that leaf holds only the variable's name. The bytes that the file will contain exist only when
the program runs, and whether an attacker could have written them decides whether this alert is
real.

Bandit reads Python only. Other SAST tools, such as Semgrep and CodeQL, do the same job for many
languages.

\subsection{A rule in Bandit's source}
\label{sec:triage-rule}

In Bandit's source, the simplest rules are short data structures. Each holds a list of function
names, a message, a CWE and a severity, and Bandit calls such a list of names a blacklist. Other
rules, such as B102 for \texttt{exec} and B608 for SQL text, are short functions. Each of them
inspects the matched node and returns a message, a CWE, a severity and a confidence.

\Cref{lst:triage-b301} shows the entry for B301, the rule that matched the line of
\cref{fig:triage-tree} \citep{bandit}. The function \texttt{build\_conf\_dict} builds the entry from
a name, an identifier, a CWE, the names to match, the message and an optional severity.

\noindent\begin{minipage}{\linewidth}
\begin{lstlisting}[breakatwhitespace=true, basicstyle=\ttfamily\footnotesize, caption={Bandit's entry for rule B301, with the list of names shortened.}, label=lst:triage-b301]
utils.build_conf_dict(
    "pickle",
    "B301",
    issue.Cwe.DESERIALIZATION_OF_UNTRUSTED_DATA,
    [
        "pickle.loads",
        "pickle.load",
        "pickle.Unpickler",
        # ... nine more names, from dill.loads
        # to pandas.read_pickle
    ],
    "Pickle and modules that wrap it can be unsafe when used to "
    "deserialize untrusted data, possible security issue.",
)
\end{lstlisting}
\end{minipage}

The parts of the entry are the following:
\bi
\item \textbf{Name and identifier.} The strings \texttt{"pickle"} and \texttt{"B301"} name the
rule. The identifier B301 appears in every alert the rule produces.
\item \textbf{CWE.} The constant \texttt{DESERIALIZATION\_OF\_UNTRUSTED\_DATA} stands for CWE-502,
deserialization of untrusted data. \textbf{Deserialization} is the step that turns stored bytes back
into objects. Every alert of the rule carries this CWE.
\item \textbf{Names to match.} The list holds the functions the rule matches, each written as a
\textbf{qualified name}: the function's name together with the module it belongs to, as in
\texttt{pickle.load}. Besides three names from \texttt{pickle}, the list holds nine names from
modules that wrap it, from \texttt{dill} to \texttt{pandas.read\_pickle}.
\item \textbf{Message.} The two strings at the end form the text reported with each alert.
\item \textbf{Severity.} The entry gives none, so its alerts take the default, MEDIUM.
\ei

The message states the reason for the rule. \texttt{pickle} is Python's format for saving objects as
bytes. Loading a pickle can make Python call functions that the bytes name, so an attacker who
controls the bytes can make the program run code of their choosing.

The rule knows nothing about what is unpickled. It matches the name of the function, so it reports
every call to \texttt{pickle.load} in the same way, whether the bytes come from a file the program
wrote itself or from a request. Every one of those alerts receives the same message, CWE and
severity.

\subsection{Severity and confidence}
\label{sec:triage-severity}

Every Bandit rule attaches two estimates to its alerts, a severity and a confidence. The
\textbf{severity} of a rule is its estimate of the harm if the flaw is real. Its
\textbf{confidence} is its estimate of how often the matched shape is that flaw. Bandit grades both
as LOW, MEDIUM or HIGH.

\Cref{tab:triage-severity} gives both estimates for the rules that fire on HW1's three targets, as
Bandit 1.9.4, the version HW1 uses, reported them. Most of these rules estimate MEDIUM harm. The
weak-hash rule B324 and the shell rule B602 estimate HIGH harm, and the rules for SQL text and for
requests without a timeout have LOW confidence.

\begin{table*}[t]
  \caption{Severity and confidence of the Bandit rules that fire on HW1's three targets.}
  \label{tab:triage-severity}
  \begin{adjustbox}{width=\linewidth}\begin{tabularx}{\linewidth}{L{0.16\linewidth}XL{0.13\linewidth}L{0.13\linewidth}}
    \toprule
    \textbf{Rule} & \textbf{What it matches} & \textbf{Severity} & \textbf{Confidence} \\
    \midrule
    \textbf{B608} & SQL text built in a string & MEDIUM & LOW \\
    \textbf{B301, B614, B615} & \texttt{pickle}, \texttt{torch.load}, downloads from the Hugging Face Hub not pinned to a commit & MEDIUM & HIGH \\
    \textbf{B102, B307} & \texttt{exec}, \texttt{eval} & MEDIUM & HIGH \\
    \textbf{B314, B317, B318} & Standard XML parsers, which are unsafe on outside data & MEDIUM & HIGH \\
    \textbf{B324} & MD5 or SHA1 in \texttt{hashlib} & HIGH & HIGH \\
    \textbf{B602} & A subprocess started through the shell & HIGH & HIGH \\
    \textbf{B113} & Requests sent without a timeout & MEDIUM & LOW \\
    \bottomrule
  \end{tabularx}\end{adjustbox}
\end{table*}

One rule adjusts its confidence to the code it matched. B608, the SQL rule, reports MEDIUM
confidence when the SQL text is built inside the call that runs it. In the development set's cases
the text is built in one statement and run in another, and there B608 reports LOW confidence.

B602, the shell rule, adjusts its severity to the code it matched in a similar way. It reports HIGH
severity when the shell command is built at run time, as in Radicale's alert, and LOW severity when
the command is a literal string. HW1's filter drops such LOW alerts.

Both estimates were made before anyone read the code around the match. They are decided when the
rule is written, and the rule assigns them from the matched shape alone. B301's entry gives no
confidence, and Bandit gives every alert of a rule built from a list of names HIGH confidence. Most
rules give every alert the same two values, and even B608 looks only at the statement it matched.

Neither estimate says where the data at the matched call comes from. B301, for example, reports HIGH
confidence on each of Radicale's pickle alerts, although whether any of them is a vulnerability
depends on who can write the files they load (the threat model of \cref{sec:discovery}).

HW1 filters alerts by severity alone. It keeps every alert of MEDIUM or HIGH severity, whatever its
confidence. A rule with LOW confidence, such as B608 or B113, therefore still produces alerts that
the agent must triage.

\subsection{One run and one alert record}
\label{sec:triage-record}

A Bandit run produces a list of alerts, and each alert reaches the agent as one structured record.
HW1's brief fixes the command for each target. For Radicale, Bandit scans the server's package
recursively, skips its tests, keeps the alerts of MEDIUM severity or higher and writes its results as
JSON. At the commit HW1 uses \citep{radicale}, this run reports ten alerts:
\bi
\item \textbf{B301 (pickle).} Seven alerts, each a call that loads pickled data, with MEDIUM
severity and HIGH confidence.
\item \textbf{B602 (shell).} One alert, a subprocess started through the shell, with HIGH severity
and HIGH confidence.
\item \textbf{B324 (weak hash).} One alert, a use of MD5 or SHA1, with HIGH severity and HIGH
confidence.
\item \textbf{B113 (timeout).} One alert, a request sent without a timeout, with MEDIUM severity and
LOW confidence.
\ei

The same command on the same commit returns the same ten alerts, as \cref{sec:discovery} explained
for every target. Any difference between the labels of two runs of a triage agent on Radicale
therefore comes from triage.

\Cref{lst:triage-alert} shows the record of the alert on the line of \cref{fig:triage-tree},
shortened. The record is JSON, and its fields name the rule, give its estimates and its CWE, and
locate the match:
\bi
\item \textbf{Rule.} \demoref{test\_id} names the rule, B301. \demoref{test\_name} is
\texttt{blacklist} for every rule built from a list of names. A rule written as a function puts the
function's name there instead, such as \texttt{hashlib} for B324.
\item \textbf{Estimates.} \demoref{issue\_severity} and \demoref{issue\_confidence} carry the rule's
two estimates.
\item \textbf{CWE.} \demoref{issue\_cwe} gives the CWE number and a link to its entry in the
catalogue. It names the kind of flaw the rule believes it found, here CWE-502.
\item \textbf{Location.} \demoref{filename} and \demoref{line\_number} say where the match is.
\ei

\noindent\begin{minipage}{\linewidth}
\begin{lstlisting}[breakatwhitespace=true, basicstyle=\ttfamily\footnotesize, caption={One Bandit alert record, as the agent receives it, with its message, code, position and documentation fields cut.}, label=lst:triage-alert]
{
  "test_id": "B301",
  "test_name": "blacklist",
  "issue_severity": "MEDIUM",
  "issue_confidence": "HIGH",
  "issue_cwe": {
    "id": 502,
    "link": "https://cwe.mitre.org/data/definitions/502.html"
  },
  "filename": "radicale/sharing/files.py",
  "line_number": 131
}
\end{lstlisting}
\end{minipage}

The full record also holds the rule's message in \demoref{issue\_text}, the exact position of the
match in \demoref{line\_range}, \demoref{col\_offset} and \demoref{end\_col\_offset}, the lines around
the match in \demoref{code}, and a link to the rule's documentation in \demoref{more\_info}.

Nothing in the record says where the loaded bytes come from. Apart from the position and the code of
the match, the record holds only what the rule supplies: its identifier, its message, its CWE and its
two estimates. Finding out where the data comes from is left to the agent.

\subsection{Vulnerability and exploit}
\label{sec:triage-vulnerability}

A weakness becomes a vulnerability when a running program lets an attacker reach it. A weakness, as
\cref{sec:discovery} defined it, is a kind of flaw, such as building a database query by pasting a
value into its text. One occurrence of a weakness in a program counts as a vulnerability only when
a threat source, such as an attacker, could exploit it.

\begin{definition}[Vulnerability]
\label{def:triage-vulnerability}
``Weakness in an information system, system security procedures, internal controls, or
implementation that could be exploited or triggered by a threat source.''
\citep{nist-glossary-vulnerability}
\end{definition}

Stated in terms of a program, a \textbf{vulnerability} is a weakness an attacker can reach in a
running program. An \textbf{exploit} is an input that reaches it. A weakness in a program is a
vulnerability when all of the following hold:
\bi
\item \textbf{A path.} The attacker's data enters at the attack surface and travels to a sink, the
call where it does harm (\cref{sec:discovery}).
\item \textbf{Nothing removes the danger.} No step on the way makes the data harmless. Binding the
value as a parameter of a query, checking the input, and putting a constant in the value's place are
examples of such steps.
\item \textbf{The harm matters.} The result is a harm that the threat model says we must prevent,
such as running code that the attacker wrote.
\ei

These conditions give a label its meaning. A label of 1 claims that a path from the attack surface
to the sink exists under the threat model, with nothing on it that removes the danger. A label of 0
claims that no such path exists.

HW1 calls an alert with a label of 1 a potential vulnerability. The word \emph{potential} marks that
triage reaches its label by reading the code, without running it. Proving the flaw with a
reproduction is the later stage of validation (\cref{sec:discovery}).

A label of 1 makes a claim about a path, so its rationale has to show that path: where the data
enters, how it travels and where it does harm. A rationale for a label of 0 names what blocks the
path, such as a step that removes the danger. The lines of code that a reviewer can read to check a
rationale are its \textbf{evidence}.

\subsection{SQL injection: a real alert and a false positive}
\label{sec:triage-sql}

The same rule can fire on a real vulnerability and on safe code, because the shape it matches is
present in both. An alert with no flaw behind it is a \textbf{false positive}, also called a false
alarm. \Crefrange{lst:triage-sql-real}{lst:triage-xpath} quote test cases of the development set,
the first 99 test cases of OWASP BenchmarkPython \citep{owasp-benchmarkpython}, whose answer key says
which cases are real. Each listing quotes the test case with its own line numbers, and a line
\texttt{...} marks lines left out.

The first case is a real SQL injection, the kind of flaw that Lecture~2's worked alert patched. In
Lecture~2, the query was built by interpolation: the code pasted a value into the query text. In
\cref{lst:triage-sql-real}, the value of a form field that the client sends is stored in a
configuration object and read back as \texttt{bar} at line 42, along the path that
\cref{sec:discovery} traced from line 31.

\noindent\begin{minipage}{\linewidth}
\begin{lstlisting}[breakatwhitespace=true, basicstyle=\ttfamily\footnotesize, caption={A real SQL injection: BenchmarkTest00099, lines 42 to 49.}, label=lst:triage-sql-real]
42  bar = conf52528.get('section52528', 'keyB-52528')
    ...
46  sql = f'SELECT username from USERS where password = \'{bar}\''
47  con = helpers.db_sqlite.get_connection()
48  cur = con.cursor()
49  cur.execute(sql)
\end{lstlisting}
\end{minipage}

Line 46 pastes \texttt{bar} into the query text with an f-string, and line 49 runs that text. The
request's value becomes part of the statement, so it can change what the statement does. A password
such as \lstinline|' OR '1'='1|, for example, closes the quoted string early and adds a condition that
every row satisfies. That input is an exploit. The label is 1, and the answer key marks the case as
a real SQL injection, CWE-89. Bandit's B608 fires on the f-string.

The second case (\cref{lst:triage-sql-false}) draws an alert from the same rule, B608, because its
query text is also written as an f-string. Line 41 sets \texttt{bar} to a constant, and the condition
at line 42 is always true, so line 43 replaces it with \texttt{param}, the value of a cookie the
client sends. This detour, like the one in \cref{sec:discovery}, hides from a scanner that
\texttt{bar} holds the client's value. The attacker's data therefore reaches the call at line 50.

\noindent\begin{minipage}{\linewidth}
\begin{lstlisting}[breakatwhitespace=true, basicstyle=\ttfamily\footnotesize, caption={A false positive from the same rule: BenchmarkTest00011, lines 41 to 50.}, label=lst:triage-sql-false]
41  bar = "This should never happen"
42  if 'should' in bar:
43      bar = param
    ...
47  sql = f'SELECT username from USERS where password = ?'
48  con = helpers.db_sqlite.get_connection()
49  cur = con.cursor()
50  cur.execute(sql, (bar,))
\end{lstlisting}
\end{minipage}

The value travels as a parameter. As in the patch of Lecture~2, the query text holds a placeholder,
the \texttt{?}, and the value is passed beside the statement as a bound parameter, the tuple
\texttt{(bar,)}. The database treats a bound parameter as data only, so the value cannot change the
statement. The f-string interpolates nothing, and B608 fires only because the string has the shape
it matches.

The label is 0, and the placeholder and the tuple are its evidence. The answer key marks this case
and its neighbor, BenchmarkTest00012, as safe. The attacker's data reaches the sink here, yet there
is no vulnerability, because a step on the way removes the danger.

\subsection{Code injection: \texttt{exec} and \texttt{eval}}
\label{sec:triage-code}

A call that runs text as Python code is dangerous only when the text can hold code an attacker
wrote. The function \texttt{exec} runs a string as Python statements, and \texttt{eval} evaluates a
string as a Python expression. Bandit reports every call to either, B102 for \texttt{exec} and B307
for \texttt{eval}.

In \cref{lst:triage-code-real}, \texttt{param} holds the value of a cookie, which the client can set
to anything. Line 47 stores it in a configuration object, line 48 reads it back unchanged as
\texttt{bar}, and line 51 passes it to \texttt{exec}.

\noindent\begin{minipage}{\linewidth}
\begin{lstlisting}[breakatwhitespace=true, basicstyle=\ttfamily\footnotesize, caption={A real code injection: BenchmarkTest00074, lines 47 to 51.}, label=lst:triage-code-real]
47  conf90091.set('section90091', 'keyB-90091', param)
48  bar = conf90091.get('section90091', 'keyB-90091')
    ...
50  try:
51      exec(bar)
\end{lstlisting}
\end{minipage}

Whatever the cookie holds runs as Python, so the label is 1. The answer key marks the case as a real
code injection, CWE-94. Bandit's B102 attaches CWE-78, OS command injection, to the same call. The
two numbers differ, and both belong to the Injection family (\cref{sec:discovery}).

\Cref{lst:triage-code-false} passes the same kind of cookie value to \texttt{eval}. Lines 41 to 43,
left out of the listing, wrap the cookie value in two fixed strings, and the slice at line 44 cuts
them off again, so \texttt{bar} is the cookie value.

\noindent\begin{minipage}{\linewidth}
\begin{lstlisting}[breakatwhitespace=true, basicstyle=\ttfamily\footnotesize, caption={A false positive on eval: BenchmarkTest00073, lines 44 to 54.}, label=lst:triage-code-false]
44  bar = string41416[4:-17]
    ...
46  if not bar.startswith('\'') or not bar.endswith('\'') or '\'' in bar[1:-1]:
        ...
50      return RESPONSE
    ...
52  try:
        ...
54          eval(bar)
\end{lstlisting}
\end{minipage}

The check at line 46 returns early unless \texttt{bar} starts and ends with a single quote and holds
no other single quote. The only text that reaches \texttt{eval} is therefore a plain string literal,
such as \lstinline|'hello'|. Evaluating a string literal produces the string and runs no code.

The label is 0, and the check at line 46 is its evidence. The answer key marks the case as safe. B307
matched the call to \texttt{eval} as it matches every such call, and the check that removes the
danger lies outside the shape it reads.

\subsection{A weak hash}
\label{sec:triage-hash}

Some labels rest on the algorithm and its purpose, with no data flow to trace. In
\cref{lst:triage-hash}, line 65 creates an MD5 hash object, line 66 feeds it the input, and line 68
takes the digest, the fixed-length value that the hash function computes from its input. Lines 69
and 70 append the digest to a password file. A later line of the case reports that a sensitive value
was hashed and stored.

\noindent\begin{minipage}{\linewidth}
\begin{lstlisting}[breakatwhitespace=true, basicstyle=\ttfamily\footnotesize, caption={A weak hash for a sensitive value: BenchmarkTest00057, lines 65 to 70.}, label=lst:triage-hash]
65  hash = hashlib.md5()
66  hash.update(input)
    ...
68  result = hash.digest()
69  f = open(f'{helpers.utils.TESTFILES_DIR}/passwordFile.txt', 'a')
70  f.write(f'hash_value={base64.b64encode(result)}\n')
\end{lstlisting}
\end{minipage}

Here the harm comes from the algorithm itself, so the call to \texttt{hashlib.md5} plays the role of
the sink. MD5 is no longer considered strong enough to protect a sensitive value such as a password.
The flaw is the choice of MD5 for this purpose, whatever value reaches it. A label of 1 on this alert
therefore claims that a secret is protected with an algorithm too weak for it, and its rationale
traces no path from the attack surface. The same call hashing a cache key would be label 0, because
a cache key protects nothing secret.

The answer key marks all five hash cases that Bandit flags as real, with CWE-328, use of a weak
hash. Bandit's B324 attaches CWE-327, use of a broken or risky cryptographic algorithm. Both belong
to the Cryptographic Failures family.

\subsection{An XML alert and an unreported XPath query}
\label{sec:triage-xpath}

An alert can fall on a harmless line while a nearby line has the shape of a flaw that no rule
reports. In \cref{lst:triage-xpath}, B314 fires on line 55, a call to an XML parser. Parsing XML that
an attacker supplies can be dangerous, which is why the rule exists. The file parsed here, however,
is one that the program ships, so no attacker data reaches the parser. The alert is therefore a false
positive, and its label is 0.

\noindent\begin{minipage}{\linewidth}
\begin{lstlisting}[breakatwhitespace=true, basicstyle=\ttfamily\footnotesize, caption={An XML alert on a fixed file and an XPath query: BenchmarkTest00013, lines 41 to 57.}, label=lst:triage-xpath]
41  bar = "alsosafe"
42  if param:
43      lst = []
44      lst.append('safe')
45      lst.append(param)
46      lst.append('moresafe')
47      lst.pop(0)
48      bar = lst[1]
    ...
54  try:
55      root = ET.parse(f'{helpers.utils.RES_DIR}/employees.xml')
56      query = f"/Employees/Employee[@emplid=\'{bar}\']"
57      nodes = elementpath.select(root, query)
\end{lstlisting}
\end{minipage}

Line 56 builds an XPath query by pasting \texttt{bar} into it. XPath is a query language that
selects parts of an XML document, and pasting a request value into an XPath query is XPath
injection, CWE-643. Bandit has no rule for this shape.

In this case the shape is harmless. Line 41 sets \texttt{bar} to the constant
\lstinline|'alsosafe'|. When \texttt{param} is not empty, lines 43 to 48 build the list
\lstinline|['safe', param, 'moresafe']|, remove its first element and take the second of those left,
so \texttt{bar} becomes the constant \lstinline|'moresafe'|. Either way, \texttt{bar} holds a
constant at line 56.

The benchmark includes safe variants of each kind of flaw it tests, and this case is one of them.
Its answer key lists it under XPath injection and marks it safe.

The real XPath cases of the development set show what discovery misses. A real XPath injection gets
no alert at all, because Bandit has no rule for it, and the development set has two such cases. A
real flaw that is left unreported, or reported and then dismissed, is a \textbf{false negative}. The
two real XPath cases are false negatives of discovery.

\subsection{What discovery leaves undecided}
\label{sec:triage-undecided}

On the development set we can measure the scanner against the answer key, and the measurement shows
what discovery leaves for triage. The 99 cases hold 34 real flaws. Bandit reports 18 alerts in 17
files, because one file carries two alerts. Each case is one file, so we count the alerts by file.

\Cref{tab:triage-undecided} sorts the outcomes. Its first two rows divide the 17 reported files, and
its third row counts the real flaws that receive no alert. The last row picks out the path-traversal
cases, whose 8 real flaws are among those 26.

\begin{table}[htbp]
  \caption{Bandit's alerts on the development set, compared with the answer key.}
  \label{tab:triage-undecided}
  \begin{adjustbox}{width=\linewidth}\begin{tabularx}{\linewidth}{XL{0.36\linewidth}}
    \toprule
    \textbf{Outcome} & \textbf{Count} \\
    \midrule
    \textbf{Reported, and real} & 8 files \\
    \textbf{Reported, and safe} & 9 files \\
    \textbf{Real, and unreported} & 26 cases \\
    \textbf{Path-traversal cases, no rule} & 21 cases, 0 alerts \\
    \bottomrule
  \end{tabularx}\end{adjustbox}
\end{table}

The \textbf{precision} of a discovery tool is the share of its alerts that are real. Its
\textbf{recall} is the share of the real flaws that it reports. These two measures summarize the
table. Counted by file, Bandit's precision on the development set is 8/17, about 0.47, and its
recall is 8/34, about 0.24.

About half of the alerts are false positives, and most of the real flaws go unreported. Of the 34
real flaws, 26 receive no alert. Some of them belong to kinds of flaw for which Bandit has no rule.
The development set's 21 path-traversal cases (\cref{sec:discovery}) receive no alert at all, and 8
of them are real.

Triage, the second stage of vulnerability management (\cref{sec:discovery}), decides for each alert
whether the match is a flaw an attacker can reach, using the code and the threat model. It works
only on the alerts that discovery hands it. It can remove false positives and so raise precision. It
cannot recover the real flaws that received no alert, so it cannot raise recall.

\subsection{The cost of each kind of mistake}
\label{sec:triage-cost}

A triage stage can be wrong in two directions, and the two mistakes cost different people. Applied to
triage, a false positive is a safe alert labeled 1 and passed on to a reviewer, and a false negative
is a real flaw labeled 0 and dismissed. For a triage stage, precision is the share of the alerts
labeled 1 that are real, and recall is the share of the real flaws among its alerts that it labels 1.
\bi
\item \textbf{A false positive passed on.} It costs a reviewer's time, because someone must read the
code to find that nothing is wrong. At scale it also costs the reviewers' trust in the tool, which is
useful only while people act on what it reports.
\item \textbf{A real flaw dismissed.} The flaw stays in the released product, and the cost falls on
the users and systems that an attacker harms through it.
\ei

A 2026 study, available as a preprint, measured both mistakes for agents that filter the alerts of
SAST tools \citep{xiong2026sifting}. It ran four tools, CodeQL, Joern, Semgrep and SonarQube, on the
Java version of the OWASP Benchmark \citep{owasp-benchmark}, a suite of test cases of the same kind
as BenchmarkPython. Together the four tools flagged 98.3\% of the benchmark's non-vulnerable cases.

Agents built on several agent frameworks then reviewed the alerts. The best configuration paired
SWE-agent with the model Claude Sonnet 4. SWE-agent is the published reactive agent that Lecture~2
compared with Agentless. In this configuration, the share of non-vulnerable cases still flagged fell
to 6.3\%.

On the benchmark's real vulnerabilities, the same configuration labeled 22.25\% as false positives,
about a fifth. These are false negatives of triage, real flaws that the agent dismissed. Such agents
are therefore suited to supporting a reviewer's decision. Letting them suppress alerts automatically
and unconditionally would have discarded about a fifth of the real flaws on this benchmark.

The study shows, on one benchmark, that an agent can remove most false positives and that the same
agent dismisses a substantial share of real flaws. It is a result on Java test cases. It does not
measure Python code, Bandit, or the agent HW1 asks for.

The high share of false positives is not limited to benchmarks. In one industry report, about 91\%
of the alerts that SAST tools raised were noise, that is, false positives (Ghost Security, 2025, as
quoted by \citet{feiglin2026sastbench}).

On the development set, comparing the labels of the agent you build in HW1 with the answer key
measures both kinds of mistake.

\subsection{Prioritization}
\label{sec:triage-priority}

When more alerts are real than a team can fix at once, triage also orders them, using the signals it
has. \textbf{Prioritization} is the ordering of alerts by expected harm, so that a reviewer's time
goes to the worst plausible flaw first. The signals differ in what they cost to obtain:
\bi
\item \textbf{Severity and confidence.} They come from the rule (\cref{sec:triage-severity}), so they
cost nothing to read. Most rules give the same pair to every alert, so these signals rarely separate
two alerts of the same rule.
\item \textbf{Harm in context and reachability.} The code shows how much harm a flaw of the alert's
family would do in this program, and whether data from an entry point can reach the flagged line. An
agent learns both by reading that code, which costs model calls.
\item \textbf{Signals about exploitation.} Public sources describe how published vulnerabilities are
exploited. The Known Exploited Vulnerabilities catalogue of the US Cybersecurity and Infrastructure
Security Agency (CISA) lists flaws that have been exploited in the wild \citep{cisa-kev}. The Exploit
Prediction Scoring System (EPSS) estimates the probability that a published vulnerability will be
exploited in the wild in the next 30 days \citep{first-epss}. Both describe published
vulnerabilities that have a CVE identifier (\cref{sec:discovery}). HW1 leaves both out.
\ei

HW1's label is the coarsest priority a triage stage can give. It sorts the alerts into two groups,
the potential vulnerabilities that deserve a reviewer's time (label 1) and the rest (label 0), with
no order inside either group.

\subsection{What a label must carry}
\label{sec:triage-label}

A label is useful only when a reviewer can check it, so each label carries a rationale and evidence.
The rationale, as \cref{sec:discovery} defined it, is the text that justifies the label. HW1 records
the evidence as a list of line ranges in the target's files. \Cref{lst:triage-row} shows one row of
HW1's output file, \texttt{report.json}, in the format of the brief, filled in for the eval case of
\cref{sec:triage-code}.

\noindent\begin{minipage}{\linewidth}
\begin{lstlisting}[breakatwhitespace=true, basicstyle=\ttfamily\footnotesize, caption={One row of report.json for the alert on eval in BenchmarkTest00073, with three fields cut.}, label=lst:triage-row]
{
  "alert_id": "owasp-BenchmarkTest00073-py-54-B307",
  "test_id": "B307",
  "line": 54,
  "label": 0,
  "status": "classified",
  "rationale": "Attacker: an anonymous client who controls the request. Line 46 admits only a quoted literal, so eval runs no code.",
  "evidence": [
    {"path": "testcode/BenchmarkTest00073.py",
     "start_line": 44, "end_line": 54}
  ]
}
\end{lstlisting}
\end{minipage}

The fields of the row are the following:
\bi
\item \textbf{Identity.} \demoref{alert\_id} names the alert by target, test case, line and rule.
\demoref{test\_id} and \demoref{line} repeat the rule and the line. The full row also carries the
alert's path, severity and confidence.
\item \textbf{Label.} The label is 0, for the reason given in \cref{sec:triage-code}.
\item \textbf{Status.} The row's status says whether the agent classified the alert. It is
\demoref{classified} when the agent did, \demoref{budget\_exhausted} when the run's token budget ran
out before the agent classified the alert, and \demoref{error} when classification failed. HW1's
token budget caps the tokens that all model calls of one run may use, across every alert. It differs
from the planning agent's call budget, which counts model calls. The row's status is the status of
this row and differs from the terminal status of a run (\cref{sec:routing}).
\item \textbf{Rationale.} Two sentences state the threat model and the reason: the check at line 46
admits only a quoted literal, so \texttt{eval} runs no code.
\item \textbf{Evidence.} Lines 44 to 54 of the test case hold the line that sets \texttt{bar}, the
check and the call to \texttt{eval}.
\ei

A reviewer accepts a label only when the rationale and the evidence let them check it. With this
row, a reviewer reads the cited lines and can confirm the label without repeating the agent's
investigation. A label of 1 needs the same support. Its rationale names the entry point, the path and
the sink (\cref{sec:triage-vulnerability}), or, for a weak hash, the secret that the algorithm
protects (\cref{sec:triage-hash}), and its evidence covers each of them.

Every rationale also states the threat model it assumed (\cref{sec:discovery}). For BenchmarkPython
that is an anonymous web client who controls the whole request, as the rationale of
\cref{lst:triage-row} states.

\subsection*{Key points}
\bi
\item Bandit parses each Python file into a syntax tree and reports an alert wherever a rule's shape
matches, such as a call to a named function, so an alert records a matched shape and nothing about
where the argument's value came from.
\item The simplest Bandit rules, such as B301, are lists of qualified function names with a CWE, a
message and a severity, and every rule assigns its severity and confidence from the matched shape
alone, so neither estimate says whether an attacker can reach a particular match.
\item A vulnerability is a weakness an attacker can reach in a running program, and a label of 1
claims, under the threat model, a path from the attack surface to the sink with nothing on it that
removes the danger, or, for a weak hash, a secret protected by too weak an algorithm.
\item The same rule fires on real vulnerabilities and on false positives, because a bound parameter,
a check on the input or a constant value removes the danger without changing the shape the rule
matches.
\item On the development set, about half of Bandit's alerts are false positives (precision about
0.47 by file) and most real flaws receive no alert (recall about 0.24), and triage can raise the
first figure but never the second.
\item A false positive passed on by triage costs reviewers time and trust, a dismissed real flaw stays
in the released product, and in one study agents removed most false positives while dismissing about
a fifth of the real vulnerabilities.
\item HW1's label of 1 or 0 is the coarsest form of prioritization, the ordering of alerts by expected
harm, and a reviewer accepts it only when its rationale and line evidence let them check it.
\ei
